% ---------------------------------------------------------------------
\section{Tartışma}\label{sec:discussion}
% ---------------------------------------------------------------------

\subsection{Yorum}\label{sec:interpretation}

% Bulguların yorumu yazarlar tarafından yazılmalıdır (SAUCIS yapay zekâ
% politikası). Yönlendirmeler, her paragrafın dayandığı analiz çıktısını
% belirtir; burada önceden yazılmış hiçbir yorum yoktur.

\TODO{Paragraf 1, yazarlar tarafından: 10 ve 20 şarj ünitesiyle eğitilen ağırlıkların 200 şarj ünitesine ve stres düzeyine aktarılıp aktarılmadığını belirtin. \Fig{fig:zeroshot}, \Tab{S-tab:t2} içindeki aktarım açığı ve yedek kural sütunları, \Tab{S-tab:t3b} içindeki bağlayıcılık oranları ve \Tab{S-tab:ts2} kullanılmalıdır}.

\TODO{Paragraf 2, yazarlar tarafından: \Fig{fig:cost} ve \Tab{S-tab:t5} içindeki garantinin maliyetini, \Tab{S-tab:t3} içindeki Attention+penalty ihlalleriyle ilişkilendirin. Birincil ölçü olarak $\Delta J^{\sm{pen}}$ ve stats/tests\_T5.csv kullanılmalıdır}.

\TODO{Paragraf 3, yazarlar tarafından: \Tab{S-tab:t3} ve \Tab{S-tab:t3b} içindeki E4'ü ve açık sınırına göre eksik enerjiyi tartışın. Her ikisini A5'in kapalı döngü niteliğiyle ilişkilendirin}.

\TODO{Paragraf 4, yazarlar tarafından: yayımlanmış EV2Gym sayılarıyla karşılaştırmayı yalnızca literatür yapılandırması üzerinde (\Tab{S-tab:ts3}) yapın; farklı simülatörler ya da senaryolar arasında asla karşılaştırma yapmayın}.

\subsection{Şarj noktası işletmecileri için pratik çıkarımlar}\label{sec:implications}

Bu alt bölüm, deneylerden sonra yazarlar tarafından aşağıdaki noktaları kapsayacak şekilde yazılacaktır.
\begin{itemize}
\item \TODO{Trafo limiti her ham eylem için sağlandığında bir CPO'nun ne kazandığı ve koşullu ayrılış sonucunun hizmet taahhüdü açısından ne anlama geldiği.}
\item \TODO{\Tab{S-tab:t5} üzerinden, belirli bir boyuttaki istasyon için hizmet verilen EV başına avro cinsinden garantinin maliyeti.}
\item \TODO{\Tab{S-tab:t7} üzerinden, sürücülerin ayrılış zamanlarını bildirmediği durumlarda konformal varyantın kullanımı ve kabul edilebilir kapsama düzeyleri.}
\item \TODO{\Tab{S-tab:t6} üzerinden, 15 dakikalık kontrol adımına göre kontrol adımı başına çıkarım süresi.}
\end{itemize}

\subsection{Sınırlılıklar}\label{sec:limitations}

Teorem~\ref{thm:invariance} kapsamındaki ayrılış sonucu, bildirilen ayrılış zamanlarına ve gereksinimlere (A2) dayanmaktadır.
Konformal genişletmemiz yalnızca ayrılış adımını gevşetmekte, kapsaması marjinal kalmakta ve son tarihlerin stratejik olarak yanlış bildirilmesi \cite{ferragut2022scheduling} modelin dışında kalmaktadır.
\Sec{sec:invariance} içindeki iki EV'li örneğin gösterdiği gibi toplam olurluluk (A5) kapalı döngüde sağlanamayabilir.
Bu tür adımlarda ayrılış gereksinimi yalnızca Teorem~\ref{thm:invariance}(b) kapsamındaki EV'ler için garanti edilir; Lemma~\ref{S-lem:deficit} ise kalan eksik enerjiyi ortadan kaldırmadan sınırlar.
Gelecek boyunca ortak olurluluğu sertifikalayan bir katman, zamana bağlı kısıtlara ihtiyaç duyacak ve tek toplamlı izdüşüm yapısını kaybedecektir.

Trafoya ilişkin sonuç, A3 tek adım modelini ve A6'daki dışsal kapasite payını gerektirir; denetleyici A6'yı uygulamaz.
Lemma~\ref{S-lem:a3}, kullandığımız EV2Gym ayarları için A3'ü ispatlamaktadır; buna karşın gerçek bir batarya yönetim sistemi komut edilenden daha az güç çekebilir.
Eksik teslimat, bir EV'nin enerjisini sınırların varsaydığı değerin altına düşürür; bu nedenle donanım uygulamasında eşikler üzerinde bir pay ya da ölçülen enerjiden geri besleme gerekecektir; bu çözüm burada analiz edilmemiştir.
A1, EV2Gym'in üstel şarj gücü azalması için geçerlidir; ancak çok daha dik bir azalma bu varsayımı bozacaktır.
Model, gerilim limitleri ya da hat akışları olmaksızın tek bir trafonun arkasındaki tek bir istasyonu kapsar ve ek bağlantılı toplam kısıtları Önerme~\ref{prop:projection} kapsamının dışında kalacaktır.

Deneylerde EV2Gym'in kayıpsız muhasebesi kullanılmakta ve deşarj fiyatı şarj fiyatına eşit alınmaktadır.
Dolayısıyla V2G arbitrajı ne gidiş-dönüş kaybı ne de yıpranma maliyeti içerir; bu nedenle deşarjdan elde edilen kârlar iyimser olabilir ve duyarlılık çalışmaları bu etkiyi yalnızca tek bir istasyon boyutunda incelemektedir.
Tüm sonuçlar, EV2Gym istatistiklerinden üretilen sentetik varışlar ve Hollanda gün öncesi fiyatlarıyla yapılan simülasyonlardan elde edilmiştir; bazı geç akşam oturumları epizot ufkuna sığmamakta ve tahmin gürültüsü sentetiktir.
Şebeke büyüklükleri $N$ ile doğrusal olarak ölçeklendiğinden test istasyonlarının şarj ünitesi başına istatistikleri eğitimdekilerle örtüşmekte; bunları yalnızca stres düzeyi değiştirmektedir.

Bazı karşılaştırmalar tam olarak eşdeğer koşullarda yapılmamıştır.
EV-GNN, kendi algoritması ve ödülüyle, tek bir boyutta, her profil için ayrı olarak, çalıştırma başına tek bir yük ve PV çekilişiyle ve varsayılanından daha az adımla eğitilmektedir.
eMPC-HiGHS yeniden gerçeklememiz ufku içindeki kesin fiyatları ve varışları bilmekte, MLP yalnızca $N\le20$ için raporlanmakta ve GNN-DT çalıştırılmamıştır.
Optimumlar, tek adımlık bir geçiş bandı ve \%1'lik bir pay dışında kesin olan dış şarj gücü azalması satırlarını yuvarlama payıyla birlikte kullanmaktadır.
Ayrıca optimumlar çevrim içi bir denetleyicinin bilgi kısıtlarını göz ardı etmektedir; dolayısıyla her açığın bir kısmı geleceği bilmemenin bedelidir.
